% Motor: XeLaTeX. Sesion 19 - material docente UTP.
% Misma estructura docente que la Sesion 18.
\documentclass[aspectratio=169,10pt,c]{beamer}
% Reutiliza sin cambios la identidad visual UTP validada en la Sesion 09.
% Reutiliza sin cambios la identidad visual UTP validada en la Sesion 08.
% Reutiliza sin cambios la identidad visual UTP validada en la Sesion 05.
% La segunda ruta permite compartir los assets comunes sin duplicar binarios.
\input{../../Sesion05/slides/header.tex}
\graphicspath{{../../Sesion05/slides/assets/}}



% El mapa de regiones es la captura del curso que ya esta en la Sesion 17.
\graphicspath{{assets/}{../../Sesion17/slides/assets/}{../../Sesion05/slides/assets/}}
% Bloques de código con el mismo estilo que la Sesion 11.
\usepackage{listings}
\lstset{basicstyle=\ttfamily\small,backgroundcolor=\color{utpGrisClaro},
  frame=none,aboveskip=4pt,belowskip=4pt,showstringspaces=false,
  columns=fullflexible,keepspaces=true,breaklines=true,
  xleftmargin=5pt,xrightmargin=5pt,framesep=5pt}


\title{Integración de Servicios Cloud GCP}
\subtitle{Ruta Associate Cloud Engineer: configuración de redes VPC}
\date{Semana 10 - Sesión 19}
\author{Rel Guzman}

% Separador de tema: titulo grande, subtitulo y linea roja.
\newcommand{\separador}[2]{%
\begin{frame}[plain]
\vfill
\begin{center}
{\fontsize{29}{34}\selectfont\bfseries #1}\par
\vspace{2mm}
{\small #2}\par
\vspace{4mm}
{\color{utpRojo}\rule{8.5cm}{1.2pt}}
\end{center}
\vfill
\end{frame}}

\begin{document}

{\setbeamertemplate{footline}{}\maketitle}

% ------------------------------------------- Revision de la sesion anterior
\begin{frame}{Revisión de la sesión anterior (1/2): una web en minutos}
\begin{block}{Lo que hicimos}
\small En el lab \emph{Implementa una pila LAMP con Google Cloud Marketplace} publicaste una web completa \textbf{sin instalar nada a mano}.
\end{block}
\begin{columns}[T,onlytextwidth]
\column{.48\textwidth}
\begin{exampleblock}{Marketplace}
\small Una \textbf{tienda} dentro de la consola con programas ya preparados.\par\smallskip
\emph{Se parece a:} la Play Store o la App Store.
\end{exampleblock}
\column{.48\textwidth}
\begin{exampleblock}{Pila LAMP}
\small Cuatro programas, uno sobre otro, para hacer un sitio web: \textbf{L}inux, \textbf{A}pache, \textbf{M}ySQL y \textbf{P}HP.\par\smallskip
\emph{Se parece a:} un restaurante con local, mozo, despensa y cocinero.
\end{exampleblock}
\end{columns}
\par\vspace{1mm}
\begin{alertblock}{La idea que seguimos usando hoy}
\small Marketplace no hizo magia: hizo \textbf{los mismos pasos} que harías a mano. Hoy aprendes a hacer tú los de la \textbf{red}.
\end{alertblock}
\end{frame}

\begin{frame}{Revisión de la sesión anterior (2/2): pregunta rápida}
\begin{center}
\begin{tikzpicture}[font=\small,
  pieza/.style={draw=utpRojo,rounded corners=3pt,fill=white,minimum width=2.75cm,minimum height=1.0cm,align=center}]
  \node[pieza] at (0,0) {\textbf{VM}\\\footnotesize\texttt{lamp-1-vm}};
  \node[pieza,fill=utpRojo!20] at (3.2,0) {\textbf{Dirección}\\\footnotesize en Internet};
  \node[pieza,fill=utpRojo!20] at (6.4,0) {\textbf{Firewall}\\\footnotesize deja pasar la web};
  \node[pieza] at (9.6,0) {\textbf{Programas}\\\footnotesize los de LAMP};
\end{tikzpicture}
\end{center}
\begin{columns}[T,onlytextwidth]
\column{.48\textwidth}
\begin{block}{Lo que Marketplace creó por ti}
\small Una máquina, una dirección para visitarla, una puerta abierta para la web y los programas instalados.
\end{block}
\column{.48\textwidth}
\begin{exampleblock}{Las dos piezas de hoy}
\small La \textbf{dirección} y el \textbf{firewall} son piezas de la \textbf{red}. Hoy las ves por dentro.
\end{exampleblock}
\end{columns}
\par\vspace{1mm}
\begin{alertblock}{Pregunta rápida}
\small Si Marketplace \textbf{no} hubiera abierto el firewall, ¿qué habría pasado al hacer clic en la \emph{URL del sitio}?
\end{alertblock}
\end{frame}

\begin{frame}[c]
\begin{block}{Logro de la sesión}
Al finalizar, podrás explicar con tus palabras \textbf{cómo se organiza Google Cloud} (la jerarquía, las regiones y las zonas), decir qué son una \textbf{red local}, una \textbf{subred} y un \textbf{gateway}, y \textbf{crear una red VPC} con sus subredes y su firewall.
\end{block}
\end{frame}

\begin{frame}{Saberes previos}
\begin{itemize}\setlength\itemsep{10pt}
  \item En tu computadora guardas archivos dentro de carpetas, y carpetas dentro de otras. ¿Para qué sirve ese orden?
  \item En tu casa, el celular, la laptop y el televisor usan el mismo wifi. ¿Pueden ``verse'' entre ellos? ¿Y con el wifi del vecino?
  \item ¿Quién decide quién entra a un edificio? ¿Cómo lo decide?
\end{itemize}
\end{frame}

\begin{frame}[c]{Utilidad de esta sesión}
\begin{center}
\begin{minipage}{.86\textwidth}
{\large Toda máquina en la nube vive \textbf{dentro de una red}. Quien entiende la red sabe responder dos preguntas de cualquier trabajo real: \textbf{¿quién se habla con quién?} y \textbf{¿quién puede entrar?}}
\end{minipage}
\end{center}
\end{frame}

\begin{frame}{Plan de hoy: teoría y un lab}
\renewcommand{\arraystretch}{1.35}
\centering\small
\begin{tabular}{@{}p{1.4cm}>{\raggedright\arraybackslash}p{6.4cm}>{\raggedright\arraybackslash}p{5.1cm}@{}}
\toprule
\textbf{Parte} & \textbf{Qué vemos} & \textbf{Del curso}\\
\midrule
Teoría & Jerarquía, regiones y zonas; red local, subred, gateway y VPC & Módulo \emph{Redes virtuales}\\
\textbf{Lab} & Crear redes y probar quién habla con quién & \emph{Crea y configura redes de VPC}\\
\bottomrule
\end{tabular}
\par\vspace{3mm}
\begin{alertblock}{Cómo lo trabajamos}
\small Vamos a lo \textbf{esencial}: una idea por diapositiva y una comparación de la vida diaria. Del lab vemos solo \textbf{lo más importante}; los pasos están en su guía.
\end{alertblock}
\end{frame}

% ------------------------------------------------ Redes VPC
\separador{Configuración de redes VPC}{De la jerarquía de Google Cloud a tu primera red}

\begin{frame}{La jerarquía de Google Cloud}
\begin{columns}[c,onlytextwidth]
\column{.52\textwidth}
\centering
\begin{tikzpicture}[font=\footnotesize,
  n/.style={draw=utpRojo,rounded corners=3pt,minimum height=.62cm,minimum width=1.75cm,align=center}]
  \node[n,fill=utpRojo!32,minimum width=2.6cm] (o) at (3.2,3.6) {\textbf{Organización}};
  \node[n,fill=utpRojo!20] (c1) at (1.6,2.4) {\textbf{Carpeta}};
  \node[n,fill=utpRojo!20] (c2) at (4.8,2.4) {\textbf{Carpeta}};
  \node[n,fill=utpRojo!10] (p1) at (0.6,1.2) {\textbf{Proyecto}};
  \node[n,fill=utpRojo!10] (p2) at (2.6,1.2) {\textbf{Proyecto}};
  \node[n,fill=utpRojo!10] (p3) at (4.8,1.2) {\textbf{Proyecto}};
  \node[n,fill=white,minimum width=.9cm] (r1) at (0.1,0) {VM};
  \node[n,fill=white,minimum width=.9cm] (r2) at (1.2,0) {Red};
  \node[n,fill=white,minimum width=1.2cm] (r3) at (2.6,0) {Bucket};
  \node[n,fill=white,minimum width=.9cm] (r4) at (4.8,0) {VM};
  \draw (o) -- (c1); \draw (o) -- (c2);
  \draw (c1) -- (p1); \draw (c1) -- (p2); \draw (c2) -- (p3);
  \draw (p1) -- (r1); \draw (p1) -- (r2); \draw (p2) -- (r3); \draw (p3) -- (r4);
  \node[anchor=west,font=\scriptsize] at (5.3,0) {recursos};
\end{tikzpicture}
\column{.45\textwidth}
\begin{block}{De arriba hacia abajo}
\small
\textbf{Organización:} toda la empresa.\par\smallskip
\textbf{Carpeta:} un área o un equipo.\par\smallskip
\textbf{Proyecto:} la caja donde se crean las cosas.\par\smallskip
\textbf{Recurso:} cada cosa que creas: una VM, un bucket, una red.
\end{block}
\end{columns}
\par\vspace{1mm}
\begin{columns}[T,onlytextwidth]
\column{.48\textwidth}
\begin{exampleblock}{Una comparación}
\small Las \textbf{carpetas de tu computadora}: el disco, las carpetas, las subcarpetas y, al final, los archivos.
\end{exampleblock}
\column{.48\textwidth}
\begin{alertblock}{Lo que se decide arriba, baja}
\small Un permiso dado en una carpeta vale para \textbf{todo lo que está dentro}. Y todo recurso pertenece a \textbf{un} proyecto.
\end{alertblock}
\end{columns}
\end{frame}

\begin{frame}{Regiones y zonas: dónde vive cada recurso}
\begin{block}{Dos preguntas distintas}
\small La \textbf{jerarquía} dice \textbf{de quién} es un recurso (su proyecto). La \textbf{región} y la \textbf{zona} dicen \textbf{dónde está} físicamente.
\end{block}
\renewcommand{\arraystretch}{1.3}
\centering\small
\begin{tabular}{@{}p{2.2cm}>{\raggedright\arraybackslash}p{4.8cm}p{2.7cm}>{\raggedright\arraybackslash}p{2.8cm}@{}}
\toprule
\textbf{Alcance} & \textbf{Qué significa} & \textbf{Ejemplo de hoy} & \textbf{Como\ldots}\\
\midrule
\textbf{Global} & Existe en todo el mundo a la vez & La red VPC & la empresa\\
\textbf{Regional} & Vive en una región: una zona geográfica & La subred & su sede en una ciudad\\
\textbf{Zonal} & Vive en una zona: un centro de datos de esa región & La máquina virtual & un edificio de esa sede\\
\bottomrule
\end{tabular}
\par\vspace{2mm}
\begin{columns}[T,onlytextwidth]
\column{.48\textwidth}
\begin{exampleblock}{Cómo leer un nombre}
\small \texttt{us-central1} es una \textbf{región}. \texttt{us-central1-a} es una \textbf{zona} de esa región: la letra final la delata.
\end{exampleblock}
\column{.48\textwidth}
\begin{alertblock}{Por eso el lab te las da}
\small Cada lab te indica \textbf{su} región y \textbf{su} zona. Si usas otras, la tarea no se pone en verde.
\end{alertblock}
\end{columns}
\end{frame}

\begin{frame}{Antes de la nube: la red local}
\begin{columns}[c,onlytextwidth]
\column{.50\textwidth}
\centering
\begin{tikzpicture}[font=\footnotesize,
  eq/.style={draw,fill=utpRojo!20,rounded corners=2pt,minimum width=1.25cm,minimum height=.6cm}]
  \draw[utpRojo,thick,rounded corners=6pt,fill=utpRojo!5] (0,0) rectangle (6.4,3.2);
  \node[anchor=north west,font=\small\bfseries] at (0.1,3.15) {Red local de una oficina};
  \node[draw=utpRojo,fill=white,rounded corners=3pt,minimum width=1.9cm,minimum height=.65cm] (sw) at (3.2,1.9) {\textbf{Switch}};
  \node[eq] (a) at (1.0,0.6) {PC};
  \node[eq] (b) at (2.5,0.6) {Laptop};
  \node[eq] (c) at (4.0,0.6) {Impresora};
  \node[eq] (d) at (5.5,0.6) {PC};
  \draw (sw) -- (a); \draw (sw) -- (b); \draw (sw) -- (c); \draw (sw) -- (d);
\end{tikzpicture}
\column{.47\textwidth}
\begin{block}{Qué es}
\small \textbf{Red local} (\emph{LAN}): varios equipos \textbf{conectados entre sí} en un mismo lugar: una casa, una oficina, un laboratorio.
\end{block}
\begin{block}{Qué se necesita}
\small Cables o wifi, un aparato que los une (el \textbf{switch}) y una \textbf{dirección IP} distinta para cada equipo.
\end{block}
\end{columns}
\par\vspace{1mm}
\begin{columns}[T,onlytextwidth]
\column{.48\textwidth}
\begin{exampleblock}{Una comparación}
\small El \textbf{wifi de tu casa}: por eso mandas un video del celular al televisor sin pasar por Internet.
\end{exampleblock}
\column{.48\textwidth}
\begin{alertblock}{La idea que nos llevamos}
\small Dentro de la misma red local, los equipos se hablan \textbf{directo}.
\end{alertblock}
\end{columns}
\end{frame}

\begin{frame}{Antes de empezar: cinco palabras de redes}
\renewcommand{\arraystretch}{1.4}
\centering\small
\begin{tabular}{@{}p{2.6cm}p{5.4cm}p{4.8cm}@{}}
\toprule
\textbf{Palabra} & \textbf{Qué es} & \textbf{Se parece a}\\
\midrule
\textbf{Red} & Computadoras conectadas entre sí & El wifi de tu casa\\
\textbf{Dirección IP} & El número que identifica a cada una & La dirección de una casa\\
\textbf{Puerto} & Un número de ``puerta'' dentro de la computadora & La ventanilla 3 del banco\\
\textbf{Protocolo} & El idioma en que conversan & Hablar en castellano o por señas\\
\textbf{Ping} & Un saludo corto para ver si la otra responde & Gritar ``¿me escuchas?''\\
\bottomrule
\end{tabular}
\par\vspace{3mm}
\begin{exampleblock}{Tres nombres que verás en el lab}
\small \textbf{SSH:} entrar a la terminal de otra máquina (puerto 22). \textbf{RDP:} ver el escritorio de una máquina Windows (puerto 3389). \textbf{ICMP:} el protocolo que usa \texttt{ping}.
\end{exampleblock}
\end{frame}

\begin{frame}{Subredes: una red partida en pedazos}
\begin{columns}[c,onlytextwidth]
\column{.50\textwidth}
\centering
\begin{tikzpicture}[font=\footnotesize,
  eq/.style={draw,fill=utpRojo!20,rounded corners=2pt,minimum width=.9cm,minimum height=.5cm}]
  \draw[utpRojo,thick,rounded corners=6pt,fill=utpRojo!5] (0,0) rectangle (6.4,3.5);
  \node[anchor=north west,font=\small\bfseries] at (0.1,3.45) {Red de la empresa};
  \draw[rounded corners=4pt,fill=white] (0.3,1.65) rectangle (6.1,2.85);
  \node[anchor=west] at (0.4,2.6) {\textbf{Subred Ventas} \texttt{10.0.1.0/24}};
  \node[eq] at (1.2,2.05) {PC}; \node[eq] at (2.4,2.05) {PC}; \node[eq] at (3.6,2.05) {PC};
  \draw[rounded corners=4pt,fill=white] (0.3,0.25) rectangle (6.1,1.45);
  \node[anchor=west] at (0.4,1.2) {\textbf{Subred Contabilidad} \texttt{10.0.2.0/24}};
  \node[eq] at (1.2,0.65) {PC}; \node[eq] at (2.4,0.65) {PC};
\end{tikzpicture}
\column{.47\textwidth}
\begin{block}{Qué es una subred}
\small Un \textbf{pedazo} de la red, con su propio \textbf{grupo de direcciones}.
\end{block}
\begin{block}{Para qué se parte la red}
\small \textbf{Orden:} cada área tiene su grupo.\par
\textbf{Seguridad:} se decide qué pasa de un grupo a otro.
\end{block}
\end{columns}
\par\vspace{1mm}
\begin{columns}[T,onlytextwidth]
\column{.48\textwidth}
\begin{exampleblock}{Una comparación}
\small Un \textbf{edificio} es la red; cada \textbf{piso} es una subred; cada \textbf{oficina} del piso es un equipo.
\end{exampleblock}
\column{.48\textwidth}
\begin{alertblock}{Cómo se lee \texttt{10.0.1.0/24}}
\small Es el \textbf{rango} de la subred: de \texttt{10.0.1.0} a \texttt{10.0.1.255}. Son 256 direcciones.
\end{alertblock}
\end{columns}
\end{frame}

\begin{frame}{Gateway: la puerta que comunica las subredes}
\begin{center}
\begin{tikzpicture}[font=\footnotesize,
  eq/.style={draw,fill=utpRojo!20,rounded corners=2pt,minimum width=.9cm,minimum height=.5cm}]
  \draw[rounded corners=4pt,fill=utpRojo!5,draw=utpRojo,thick] (0,0) rectangle (4.3,1.7);
  \node[anchor=north west,font=\small\bfseries] at (0.1,1.65) {Subred Ventas};
  \node[eq] (a) at (1.0,0.55) {PC A}; \node[eq] at (2.4,0.55) {PC};
  \draw[rounded corners=4pt,fill=utpRojo!5,draw=utpRojo,thick] (8.9,0) rectangle (13.2,1.7);
  \node[anchor=north west,font=\small\bfseries] at (9.0,1.65) {Subred Contabilidad};
  \node[eq] at (10.8,0.55) {PC}; \node[eq] (b) at (12.2,0.55) {PC B};
  \node[draw=utpRojo,fill=utpRojo!32,rounded corners=4pt,minimum width=2.4cm,minimum height=1.0cm,align=center] (g) at (6.6,0.85) {\textbf{Gateway}\\\scriptsize (puerta de enlace)};
  \draw[latex-latex,thick] (4.3,0.85) -- (g);
  \draw[latex-latex,thick] (g) -- (8.9,0.85);
  \node[draw,rounded corners=3pt,fill=white,minimum width=1.8cm,minimum height=.55cm] (i) at (6.6,2.3) {Internet};
  \draw[latex-latex,thick] (g) -- (i);
\end{tikzpicture}
\end{center}
\begin{columns}[T,onlytextwidth]
\column{.48\textwidth}
\begin{block}{Qué es}
\small El equipo por el que pasa todo lo que \textbf{sale de una subred}: hacia otra subred o hacia Internet. En casa y en la oficina lo hace el \textbf{router}.
\end{block}
\column{.48\textwidth}
\begin{exampleblock}{Una comparación}
\small La \textbf{puerta del piso}. Para ir de Ventas a Contabilidad no atraviesas la pared: sales por tu puerta y entras por la otra.
\end{exampleblock}
\end{columns}
\par\vspace{1mm}
\begin{alertblock}{La regla}
\small ¿El destino está en \textbf{mi} subred? Le hablo directo. ¿Está \textbf{fuera}? Se lo entrego al gateway y él lo encamina.
\end{alertblock}
\end{frame}

\begin{frame}{De la oficina a la nube: Google Cloud VPC}
\begin{block}{Definición}
\small \textbf{VPC} (\emph{Virtual Private Cloud}, nube privada virtual): la \textbf{red privada} de tu proyecto dentro de Google Cloud. Es la misma idea de la red local, pero \textbf{sin cables}: se arma con clics o comandos.
\end{block}
\renewcommand{\arraystretch}{1.3}
\centering\small
\begin{tabular}{@{}p{3.0cm}p{4.5cm}p{5.4cm}@{}}
\toprule
 & \textbf{En la oficina} & \textbf{En Google Cloud}\\
\midrule
\textbf{La red} & Cables, switch y router & Una \textbf{red VPC}: se crea en segundos\\
\textbf{La subred} & Un piso o un área & Un pedazo de la VPC en \textbf{una región}\\
\textbf{El gateway} & El router que compras e instalas & Google lo pone \textbf{solo} en cada subred\\
\textbf{Los equipos} & PCs e impresoras & Máquinas virtuales\\
\bottomrule
\end{tabular}
\par\vspace{2mm}
\begin{columns}[T,onlytextwidth]
\column{.48\textwidth}
\begin{exampleblock}{El gateway, sin trabajo}
\small Dos subredes de la \textbf{misma VPC} ya se comunican: Google hace de gateway entre ellas. Tú no configuras nada.
\end{exampleblock}
\column{.48\textwidth}
\begin{alertblock}{Lo mismo, fuera de Google}
\small \textbf{VPC} en AWS y \textbf{Virtual Network} en Azure. La idea es la misma en todas las nubes.
\end{alertblock}
\end{columns}
\end{frame}

\begin{frame}{Proyecto, red y subred}
\begin{columns}[c,onlytextwidth]
\column{.54\textwidth}
\centering
\begin{tikzpicture}[font=\footnotesize]
  \draw[black!55,rounded corners=6pt] (-0.25,-0.25) rectangle (6.85,4.2);
  \node[anchor=north west,font=\small\bfseries] at (-0.15,4.15) {Proyecto};
  \draw[utpRojo,thick,rounded corners=6pt,fill=utpRojo!5] (0,0) rectangle (6.6,3.5);
  \node[anchor=north west,font=\small\bfseries] at (0.1,3.45) {Red (es mundial)};
  \draw[rounded corners=4pt,fill=white] (0.3,0.3) rectangle (3.1,2.7);
  \node[anchor=north west] at (0.35,2.65) {\textbf{Subred}};
  \node[anchor=north west] at (0.35,2.25) {región de EE.\,UU.};
  \node[draw,fill=utpRojo!20,rounded corners=2pt] at (1.05,0.9) {VM};
  \node[draw,fill=utpRojo!20,rounded corners=2pt] at (2.3,0.9) {VM};
  \draw[rounded corners=4pt,fill=white] (3.5,0.3) rectangle (6.3,2.7);
  \node[anchor=north west] at (3.55,2.65) {\textbf{Subred}};
  \node[anchor=north west] at (3.55,2.25) {región de Europa};
  \node[draw,fill=utpRojo!20,rounded corners=2pt] at (4.9,0.9) {VM};
\end{tikzpicture}
\column{.43\textwidth}
\begin{block}{De afuera hacia adentro}
\small
\textbf{Proyecto:} la caja donde vive todo.\par\smallskip
\textbf{Red:} abarca \textbf{todo el mundo}.\par\smallskip
\textbf{Subred:} un pedazo de la red en \textbf{una región}, con su propio grupo de direcciones.\par\smallskip
\textbf{VM:} vive dentro de una subred.
\end{block}
\end{columns}
\par\vspace{1mm}
\begin{columns}[T,onlytextwidth]
\column{.48\textwidth}
\begin{exampleblock}{Una comparación}
\small La red es una \textbf{empresa}; cada subred, su \textbf{oficina} en una ciudad; cada VM, un \textbf{empleado} con su anexo.
\end{exampleblock}
\column{.48\textwidth}
\begin{alertblock}{Sin red no hay VM}
\small Toda máquina virtual necesita una red. Lo comprobarás en el lab.
\end{alertblock}
\end{columns}
\end{frame}

\begin{frame}{Tres tipos de red}
\renewcommand{\arraystretch}{1.35}
\centering\small
\begin{tabular}{@{}>{\raggedright\arraybackslash}p{3.0cm}>{\raggedright\arraybackslash}p{5.6cm}p{4.2cm}@{}}
\toprule
\textbf{Tipo} & \textbf{Cómo son sus subredes} & \textbf{Como\ldots}\\
\midrule
\textbf{Predeterminada} (\texttt{default}) & Ya viene en cada proyecto, con una subred en cada región & el chip que ya trae el celular\\
\textbf{Modo automático} & Google crea una subred en cada región por ti & pedir el combo\\
\textbf{Modo personalizado} & Tú decides qué subredes, dónde y con qué direcciones & pedir a la carta\\
\bottomrule
\end{tabular}
\par\vspace{3mm}
\begin{columns}[T,onlytextwidth]
\column{.48\textwidth}
\begin{exampleblock}{¿Cuál conviene?}
\small \textbf{Automático:} para probar rápido.\par
\textbf{Personalizado:} para trabajos reales; tú tienes el control.
\end{exampleblock}
\column{.48\textwidth}
\begin{alertblock}{Un solo sentido}
\small De automático a personalizado \textbf{sí} se puede cambiar. De regreso, \textbf{no}.
\end{alertblock}
\end{columns}
\end{frame}

\begin{frame}{La idea más importante: ¿misma red u otra red?}
\begin{center}
\begin{tikzpicture}[font=\footnotesize,
  vm/.style={draw,fill=utpRojo!20,rounded corners=2pt,minimum width=1.0cm,minimum height=.6cm}]
  \draw[utpRojo,thick,rounded corners=6pt,fill=utpRojo!5] (0,0) rectangle (6.0,2.2);
  \node[anchor=north west,font=\small\bfseries] at (0.1,2.15) {Red 1};
  \node[vm] (a) at (1.2,0.75) {VM A};
  \node[vm] (b) at (4.8,0.75) {VM B};
  \node[anchor=north] at (1.2,0.4) {EE.\,UU.};
  \node[anchor=north] at (4.8,0.4) {Europa};
  \draw[latex-latex,thick] (a) -- node[above] {se hablan directo} (b);
  \draw[utpRojo,thick,rounded corners=6pt,fill=utpRojo!5] (7.0,0) rectangle (9.6,2.2);
  \node[anchor=north west,font=\small\bfseries] at (7.1,2.15) {Red 2};
  \node[vm] (c) at (8.3,0.75) {VM C};
  \draw[utpRojo,thick,rounded corners=6pt,fill=utpRojo!5] (10.4,0) rectangle (13.0,2.2);
  \node[anchor=north west,font=\small\bfseries] at (10.5,2.15) {Red 3};
  \node[vm] (d) at (11.7,0.75) {VM D};
  \node[anchor=north] at (10.0,0.4) {misma ciudad};
  \draw[dashed,thick] (c) -- node[above,font=\small\bfseries] {X} (d);
\end{tikzpicture}
\end{center}
\begin{columns}[T,onlytextwidth]
\column{.48\textwidth}
\begin{block}{A y B: misma red, países distintos}
\small Se hablan \textbf{directo}, por dentro, como si estuvieran en el mismo cuarto.
\end{block}
\column{.48\textwidth}
\begin{block}{C y D: redes distintas, misma ciudad}
\small \textbf{No} se hablan por dentro. Solo si cada una sale con su dirección pública.
\end{block}
\end{columns}
\par\vspace{1mm}
\begin{exampleblock}{Una comparación}
\small Dos anexos de la \textbf{misma empresa} se llaman gratis aunque estén en ciudades distintas. Dos \textbf{empresas distintas}, aunque compartan edificio, tienen que llamarse por el número público.
\end{exampleblock}
\end{frame}

\begin{frame}{Dirección IP interna y externa}
\begin{columns}[T,onlytextwidth]
\column{.48\textwidth}
\begin{block}{IP interna}
\small La tiene \textbf{toda} máquina virtual. Sirve para hablar \textbf{dentro} de su red. Google la asigna sola.
\end{block}
\begin{block}{IP externa}
\small Es \textbf{opcional}. Sirve para que te vean desde \textbf{Internet}. La \emph{URL del sitio} de la sesión anterior era esto.
\end{block}
\column{.48\textwidth}
\begin{exampleblock}{Una comparación}
\small La interna es tu \textbf{anexo}: solo funciona dentro de la oficina. La externa es el \textbf{número público}: cualquiera puede marcarlo.
\end{exampleblock}
\begin{alertblock}{La externa puede cambiar}
\small \textbf{Efímera:} si apagas y prendes la VM, te dan otra.\par
\textbf{Estática:} la reservas y es siempre la misma.
\end{alertblock}
\end{columns}
\par\vspace{1mm}
\begin{block}{Lo que muestra la demostración del curso}
\small Apagan una VM y la vuelven a prender: la IP \textbf{interna quedó igual}; la \textbf{externa cambió}.
\end{block}
\end{frame}

\begin{frame}{Firewall: quién puede pasar}
\begin{center}
\begin{tikzpicture}[font=\small,
  caja/.style={draw=utpRojo,rounded corners=4pt,minimum width=2.6cm,minimum height=1.0cm,align=center}]
  \node[caja,fill=white] (i) at (0,0) {\textbf{Alguien}\\\footnotesize quiere entrar};
  \node[caja,fill=utpRojo!30] (f) at (5.0,0) {\textbf{Firewall}\\\footnotesize revisa su lista};
  \node[caja,fill=utpRojo!10] (v) at (10.0,0) {\textbf{Tu VM}};
  \draw[-latex,thick] (i) -- (f);
  \draw[-latex,thick] (f) -- node[above,font=\footnotesize] {si está permitido} (v);
\end{tikzpicture}
\end{center}
\begin{columns}[T,onlytextwidth]
\column{.48\textwidth}
\begin{block}{Las dos reglas que siempre están}
\small Lo que \textbf{entra}: todo \textbf{bloqueado}.\par
Lo que \textbf{sale}: todo \textbf{permitido}.\par\smallskip
Tú agregas reglas para \textbf{dejar entrar} lo que necesitas.
\end{block}
\column{.48\textwidth}
\begin{exampleblock}{Una regla dice tres cosas}
\small
\textbf{Qué:} protocolo y puerto (SSH, puerto 22).\par
\textbf{Desde dónde:} qué direcciones.\par
\textbf{A quién:} todas las VMs o solo algunas.
\end{exampleblock}
\end{columns}
\par\vspace{1mm}
\begin{alertblock}{Una comparación}
\small El \textbf{portero} del condominio con su lista de invitados: si no estás en la lista, no entras. Tú, que vives ahí, sales cuando quieres.
\end{alertblock}
\end{frame}

% ------------------------------------------------------------------ Lab
\separador{El lab: lo más importante}{\emph{Crea y configura redes de VPC}}

\begin{frame}{El lab en una mirada}
\begin{block}{De qué trata}
\small Creas \textbf{tres redes} propias y pruebas con \texttt{ping} \textbf{qué máquinas se pueden hablar}.
\end{block}
\begin{center}
\begin{tikzpicture}[font=\footnotesize,
  vm/.style={draw,fill=utpRojo!20,rounded corners=2pt,minimum height=.55cm,inner xsep=3pt}]
  \draw[utpRojo,thick,rounded corners=6pt,fill=utpRojo!5] (0,0) rectangle (5.0,3.0);
  \node[anchor=north west,font=\small\bfseries] at (0.1,2.95) {\texttt{mynetwork}};
  \node[vm] at (2.5,1.75) {\texttt{mynet-us-vm} (EE.\,UU.)};
  \node[vm] at (2.5,0.75) {\texttt{mynet-notus-vm} (otra región)};
  \draw[utpRojo,thick,rounded corners=6pt,fill=utpRojo!5] (5.4,0) rectangle (9.5,3.0);
  \node[anchor=north west,font=\small\bfseries] at (5.5,2.95) {\texttt{managementnet}};
  \node[vm] at (7.45,1.25) {\texttt{managementnet-us-vm}};
  \draw[utpRojo,thick,rounded corners=6pt,fill=utpRojo!5] (9.9,0) rectangle (13.4,3.0);
  \node[anchor=north west,font=\small\bfseries] at (10.0,2.95) {\texttt{privatenet}};
  \node[vm] at (11.65,1.25) {\texttt{privatenet-us-vm}};
\end{tikzpicture}
\end{center}
\begin{columns}[T,onlytextwidth]
\column{.48\textwidth}
\begin{block}{Cómo leer el dibujo}
\small Tres \textbf{redes} (los marcos) y cuatro \textbf{máquinas} (las cajas). Tres máquinas están en la \textbf{misma zona}, pero en redes distintas.
\end{block}
\column{.48\textwidth}
\begin{alertblock}{La pregunta del lab}
\small Desde \texttt{mynet-us-vm}, ¿a cuáles de las otras tres llega un \texttt{ping}?
\end{alertblock}
\end{columns}
\end{frame}

\begin{frame}{Lo más importante del lab: ¿quién se puede hablar?}
\begin{block}{El experimento}
\small Entras a \texttt{mynet-us-vm} y le haces \texttt{ping} a las otras tres máquinas: primero a su IP \textbf{externa}, luego a su IP \textbf{interna}.
\end{block}
\renewcommand{\arraystretch}{1.3}
\centering\small
\begin{tabular}{@{}p{4.3cm}p{3.6cm}p{2.3cm}p{2.3cm}@{}}
\toprule
\textbf{Destino} & \textbf{Dónde está} & \textbf{IP externa} & \textbf{IP interna}\\
\midrule
\texttt{mynet-notus-vm} & misma red, otra región & Responde & \textbf{Responde}\\
\texttt{managementnet-us-vm} & otra red, misma zona & Responde & \textbf{No llega}\\
\texttt{privatenet-us-vm} & otra red, misma zona & Responde & \textbf{No llega}\\
\bottomrule
\end{tabular}
\par\vspace{2mm}
\begin{columns}[T,onlytextwidth]
\column{.48\textwidth}
\begin{exampleblock}{Qué significa}
\small Por dentro, lo que importa es estar en la \textbf{misma red}. La distancia no importa.
\end{exampleblock}
\column{.48\textwidth}
\begin{alertblock}{La conclusión del lab}
\small Cada red VPC está \textbf{aislada} de las demás. Para unirlas hay que configurarlo aparte.
\end{alertblock}
\end{columns}
\par\vspace{1mm}
{\raggedright\small \textbf{¿Por qué la IP externa sí responde?} Porque va por Internet, y tus reglas de firewall dejan entrar el \texttt{ping}.\par}
\end{frame}

\begin{frame}{El lab: tres ideas para no perderte}
\renewcommand{\arraystretch}{1.45}
\centering\small
\begin{tabular}{@{}>{\raggedright\arraybackslash}p{3.6cm}>{\raggedright\arraybackslash}p{5.3cm}>{\raggedright\arraybackslash}p{4.0cm}@{}}
\toprule
\textbf{La idea} & \textbf{Dónde la ves en el lab} & \textbf{Se parece a}\\
\midrule
\textbf{Sin red no hay VM} & Borras la red y ya no te deja crear máquinas & Pedir un anexo donde no hay central telefónica\\
\textbf{Para entrar hace falta permiso} & Una regla de firewall deja abrir la terminal solo a las VMs con la etiqueta \texttt{iap-gce} & Un sticker de ``visita autorizada''\\
\textbf{El \texttt{ping} se lee fácil} & \texttt{0\% packet loss}: respondió. \texttt{100\% packet loss}: no llegó & Gritar ``¿me escuchas?''\\
\bottomrule
\end{tabular}
\par\vspace{3mm}
\begin{columns}[T,onlytextwidth]
\column{.48\textwidth}
\begin{exampleblock}{Lo haces de dos maneras}
\small Una red con \textbf{clics} y otra con \textbf{comandos}. Hacen lo mismo: son las dos puertas de la Sesión 17.
\end{exampleblock}
\column{.48\textwidth}
\begin{alertblock}{Si algo falla}
\small Copia los nombres, la región y la zona \textbf{tal cual} los da el lab. Y sigue la guía del lab paso a paso: ahí está cada clic.
\end{alertblock}
\end{columns}
\end{frame}

% ------------------------------------------------------------ Practica
\separador{Práctica}{El lab real}

\begin{frame}{Práctica: redes VPC \textnormal{\small(en clase + tarea)}}
\begin{block}{En clase}
\small
1. Antes de empezar, responde en tu cuaderno la \textbf{pregunta del lab}: desde \texttt{mynet-us-vm}, ¿a cuáles de las otras tres llega un \texttt{ping}?\par\smallskip
2. Haz el lab real: \emph{Crea y configura redes de VPC}.\par\smallskip
3. Al llegar a la última tarea, compara lo que salió con lo que respondiste.
\end{block}
\begin{columns}[T,onlytextwidth]
\column{.48\textwidth}
\begin{exampleblock}{De tarea}
\small Terminar el lab si no alcanzó el tiempo en clase.
\end{exampleblock}
\column{.48\textwidth}
\begin{alertblock}{Entregable}
\small Captura del lab con sus tareas en verde.
\end{alertblock}
\end{columns}
\end{frame}

\begin{frame}{Cierre: preguntas para comprobar}
\begin{block}{Responde con tus palabras}
\small
1. Ordena de mayor a menor: proyecto, recurso, organización, carpeta.\par\medskip
2. Una red VPC, una subred y una máquina virtual: ¿cuál es global, cuál es regional y cuál es zonal?\par\medskip
3. ¿Para qué sirve un gateway? ¿Quién hace ese trabajo dentro de una VPC?\par\medskip
4. Dos VMs están en la misma red, una en EE.\,UU. y otra en Europa. ¿Se hablan por su IP interna? ¿Y dos VMs de redes distintas en la misma zona?\par\medskip
5. ¿Qué hace el firewall con lo que entra y con lo que sale, si no agregas ninguna regla?
\end{block}
\begin{alertblock}{Lo que viene}
\small Máquinas \textbf{sin dirección pública}: cómo llegan a Google y a Internet con \textbf{Acceso privado a Google} y \textbf{Cloud NAT}.
\end{alertblock}
\end{frame}

\begin{frame}[plain]
\begin{tikzpicture}[remember picture,overlay]
  \node[anchor=center,inner sep=0pt] at (current page.center)
    {\includegraphics[width=\paperwidth,height=\paperheight]{gracias-final.pdf}};
\end{tikzpicture}
\end{frame}

\begin{frame}{Referencias}
\scriptsize
\begin{itemize}\setlength\itemsep{3pt}
  \item Google Skills: ruta \emph{Certificación de Associate Cloud Engineer}, \texttt{skills.google/paths/11}.
  \item Google Skills: \emph{Infraestructura esencial de Google Cloud: conceptos básicos}, \texttt{skills.google/course\_templates/50}, módulo \emph{Redes virtuales}.
  \item Google Skills: lab \emph{Crea y configura redes de VPC}.
  \item Google Cloud: \emph{Resource hierarchy}, \texttt{cloud.google.com/resource-manager/docs}.
  \item Google Cloud: \emph{VPC networks} y \emph{VPC firewall rules}, \texttt{cloud.google.com/vpc/docs}.
  \item Material de repaso: \texttt{Sesion18/slides/slides.tex}.
  \item Las respuestas de todas las preguntas de esta sesión están en \texttt{respuestas.md}.
\end{itemize}
\end{frame}

\end{document}
